\section{Loading and checking the dataset}
\label{part:data}

The dataset was loaded with the MoleculeNet loader of PyTorch Geometric \parencite{PyG}, which returns each molecule as a ready-made graph. Every graph was then turned into a record holding its SMILES string (the text form of a structure), its attribute matrices and its label, and the whole set was checked before any modelling. \Cref{tab:audit} lists what the check found.

\begin{table}[htb]
    \centering
    \caption{What the audit found, and how large the molecules are. Bonds are counted once, not once per direction.}
    \label{tab:audit}
    \footnotesize
    \setlength{\tabcolsep}{3pt}
    \begin{tabular}{@{}lr@{\hspace{0.5em}}lrr@{}}
        \toprule
        Check & Value & Size & Atoms & Bonds \\
        \midrule
        Loaded / in file      & 2039 / 2050 & Mean   & 24.06 & 25.95 \\
        Unparsable, dropped   & 11         & Std    & 10.59 & 11.72 \\
        Non-finite labels     & 0          & Min    & 2     & 1     \\
        Class 0 / class 1     & 479 / 1560 & Q1     & 18    & 19    \\
        Unique canonical      & 1975       & Median & 23    & 25    \\
        Duplicated molecules  & 60         & Q3     & 28.5  & 31    \\
        \quad consistent labels & 50       & Max    & 132   & 145   \\
        \quad conflicting labels & 10      &        &       &       \\
        \bottomrule
    \end{tabular}
\end{table}

The source file holds 2050 rows but the loader returns 2039 graphs. The 11 missing rows are exactly those whose SMILES strings RDKit \parencite{RDKit} cannot read, so they are dropped before the data reaches any model. Nothing else is broken: no label is missing, no graph is empty, and every remaining structure parses.

Duplicates need more care. The same molecule can be written as several different SMILES strings, so comparing the strings directly is not enough --- and indeed all 2039 strings in the file are different, which makes it look as though there are no duplicates at all. Each structure was therefore re-written in its \emph{canonical} form, the one spelling RDKit always produces for a given molecule, and duplicates were found by comparing those. This reveals 60 molecules that appear more than once, spread over 124 rows.

Counting how many different labels each duplicated molecule carries separates two very different cases. Fifty of them repeat with the same label, which is harmless. The other ten appear with \emph{both} labels, which means the file contradicts itself about them. They are not obscure compounds: they include aspirin, levodopa and atropine, and all ten are listed in \cref{tab:conflicts}. There is no way to know which label is right, so keeping either copy would be inventing an answer; all ten molecules were removed. The remaining duplicates were reduced to one copy each. This leaves \textbf{1965 unique molecules} for the rest of the work.

The contradictions are worth taking seriously. They suggest the file was assembled from more than one source, each measuring permeability in its own way, rather than from a single consistent experiment. Whatever accuracy a model reaches here is limited by that noise as much as by the model itself.
